\chapter{Contesto e obbiettivi del progetto}
Qualcosa tipo la relazione di tirocinio sul perché sto facendo questa cosa. (guarda robe commentate).
%Il tirocinio viene svolto in forma remota per conto dell'Istituto Nazionale di Geofisica e Vulcanologia. Il tirocinio consta di un progetto di sviluppo riguardante l'infrastruttura a supporto della ricerca presso tale istituto.\\
%L'obbiettivo del progetto è quello di implementare un cluster HPC basato su Slurm funzionante sulla piattaforma OpenStack. Questo con l'obbiettivo di verificare la fattibilità dell'utilizzo di un applicativo simile in questo contesto e per automatizzare il processo di creazione dei cluster e quindi velocizzare il lavoro dei ricercatori. \\
%Il cluster viene dispiegato in un Virtual Data Center, uno spazio cloud virtuale pensato per evitare che gli utenti che lo utilizzano (come i Data Scientists) si debbano preoccupare della gestione dell'infrastruttura, fornendo uno spazio in cui hanno libertà di fare operazioni complesse. \\
%In un VDC si possono creare macchine virtuali ed altri componenti delegando agli amministratori di rete il supporto e la gestione dell'infrastruttura. 
%Il VDC permette l'accesso alle macchine via SSH ed attraverso l'interfaccia Web di OpenStack, denominata \textit{Horizon}. \\
%La possibilità di poter automatizzare il provisioning dei cluster HPC solleverebbe i ricercatori dal peso di dover conoscere e controllare l'ambiente in cui operano inoltre gli eventuali problemi potrebbero essere risolti in maniera centralizzata senza avere configurazioni sparse nella rete. \\
%
%In questo documenti si affronterà l'analisi dei componenti necessari e delle possibili alternative disponibili, l'architettura ed il funzionamento pratico del cluster. 
%
\chapter{Architettura e componenti utilizzati}
\section{OpenStack}
OpenStack è una piattaforma open-source per la gestione delle infrastrutture cloud. Nell'contesto dell'INGV è stata adottata come piattaforma per la creazione di
macchine virtuali permettendo di dividere le risorse tra diversi gruppi di utenti (detti \textit{tenant}) che possono essere i ricercatori.
Questa piattaforma è particolarmente flessibile perché, essendo divisa in componenti, permette di scegliere quali utilizzare e come configurarli. Essa può essere
una scelta valida per le imprese o istituzioni che non vogliono dipendere da fornitori di servizi esterni o hanno bisogno di soluzioni personalizzate.
Vediamo elencati ora alcuni dei principali componenti di OpenStack:
\\
\textbf{Aggiungere qualche spiegazione migliore su questi componenti (slides, documentazione)}
\\
\begin{itemize}
    \item \textbf{Nova}: è il servizio che gestisce il \textit{provisioning} delle macchine virtuali e fisiche. 
    \item \textbf{Neutron}: gestisce tutti gli aspetti della rete, permettendo la comunicazione tra risorse e componenti di OpenStack. 
    \item \textbf{Cinder}: espone l'API per la gestione dello storage a blocchi, fornendo agli utenti ed alle macchine i volumi necessari.
    \item \textbf{Glance}: mantiene le immagini e metadati che l'utente utilizza in altri servizi. 
    \item \textbf{Keystone}: fornisce autenticazione per l'intera piattaforma. Gestisce utenti, ruoli, progetti e policy di accesso.
    \item \textbf{Horizon}: l'interfaccia web che permette agli utenti di interagire con OpenStack attraverso un browser. Fornisce un'interfaccia grafica per la gestione delle risorse.
\end{itemize}
Questo progetto verrà sviluppato all'interno di un \textit{Virtual Data Center} (VDC), un ambiente virtuale creato per isolare le risorse di produzione
da quelle di sviluppo. L'idea è quella di fornire agli utenti un ambiente isolato su cui hanno controllo totale e quindi la possibilità di creare macchine virtuali,
reti, VPN, firewall ed esporre i servizi su un indirizzo IP pubblico senza preoccuparsi della gestione dell'infrastruttura sottostante e danneggiare le risorse
in produzione.
%
%\begin{figure}[H]
%    \centering
%    \includegraphics[width=0.8\textwidth]{res/vdc-diagram.png}
%    \caption{Schema di un Virtual Data Center}
%    \label{fig:vdc}
%\end{figure} 
%
\\
\\
\textbf{Domanda: è opportuno spiegare OpenStack all'interno del capitolo sul "contesto" oppure ha senso spiegarlo
qui dove si spiegano tutti i componenti?\\
Per ora lo spiego qui al massimo lo spostiamo.}
\\ 
\\
Ora che si ha compreso OpenStack potrebbe avere senso compararla con Kubernetes; questo risulta difficile perché
le due piattaforme hanno obbiettivi e contesti di applicazione diversi. Vediamo una breve panoramica delle due
messe a confronto:
\begin{itemize}
    \item \textbf{OpenStack} gestisce l'infrastruttura cloud e le risorse virtualizzate, viene implementato
        direttamente su macchine fisiche.
    \item \textbf{Kubernetes} si concentra sulla gestione ed orchestrazione di container, unità più piccole
        rispetto alle macchine virtuali ed il suo obbiettivo finale è l'esecuzione e lo scaling di applicazioni.
\end{itemize}
\textbf{Rivedere questo pippone filosofico di confrono} 
\\ 
Anche l'architettura è diversa: la prima è formata da diversi componenti interconnessi mentre la seconda è formata 
da nodi worker e nodi master che lavorano all'esecuzione di container contenuti in pod (l'unità base di Kubernets, v. sezione \ref{sec:kube}).\\
Si può dire che OpenStack è più completa e meno astratta perché punta al controllo di servizi cloud basici mentre Kubernetes lavora 
a livello applicativo, erogando direttamente il servizio usato dagli utenti.\\
Entrambe possono essere disposte su macchine fisiche ma la tendenza del momento è disporre Kubernetes già su piattaforme virtualizzate, infatti 
i principali fornitori di servizi cloud lo erogano come \textit{Platform-as-a-Service} (PaaS) pronto all'uso permettendo agli sviluppatori
di pensare a livello applicativo facilitando lo scaling verticale senza rapportarsi con macchine fisiche.\\
In conclusione le due piattaforme vengono scelte da enti con necessità di astrazione ed implementazione: chi avrà bisogno della gestione 
dell'intera infrastruttura a partire dalle macchine fisiche (OpenStack) o chi avrà bisogno di disporre ed erogare applicazioni scalabili.
\section{Cloud-Init}
Una tecnologia di \textit{Infrastructure-as-Code} (IaC) con cui le istanze vengono configurate automaticamente al loro avvio installando pacchetti, scrivendo file di configurazione e inviando comandi alla shell. \\
IaC è un approccio che automatizza la creazione e gestione dei componenti dell'infrastruttura. Favorisce la standardizzazione ed il versioning velocizzando notevolmente il lavoro degli sviluppatori.
Questo componente è già nativo in OpenStack e quindi rende lo sviluppo veloce soprattutto in questo momento che si è nelle fasi iniziali del progetto. Potrebbe non essere la soluzione giusta in un contesto di produzione infatti vi sono applicativi più completi per l'automazione del deployment di cluster HPC come:
\begin{itemize}
    \item \textbf{Heat}: il servizio integrato della suite OpenStack per orchestrare tramite template, che in questo caso non era disponibile.
    \item \textbf{Warewulf}: una piattaforma creata appositamente per il automatizzare il provisioning di cluster HPC dotato di gestione delle immagini e creazione automatica dei nodi. Molto più fluido rispetto a Heat ma non adatto per questo stato del progetto anche perché richiede una configurazione più importante delle macchine virtuali (N.B.: WareWulf esegue il provisioning mentre Heat, Ansible e Cloud-Init eseguono la configurazione della macchina quindi W.W. andrebbe utilizzato assieme a qualche altro tool di configurazione).
    \item \textbf{Ansible}: un'altra valida alternativa a Cloud-Init che permette molta più libertà ed anche la possibilità di un mantenimento continuo delle istanze. A sua volta è più complicato da configurare e richiede un ulteriore nodo master per inviare i moduli ai nodi che controlla.
\end{itemize}
%
\section{Kubernetes}
\label{sec:kube}
Kubernetes è un sistema open-source per il deployment e lo scaling di applicazioni containerizzate. E' un sistema PaaS (Platform-as-a-Service) e lavorando a livello dei container fornisce diversi tool come scaling e load balancing ma lascia agli utenti l'integrazione di logging e monitoring. Ciò apre ad una particolare flessibilità e possibilità di integrazione con parecchie piattaforme cloud, infatti Kubernetes è fornito come servizio \textit{managed} dai principali fornitori cloud tra cui: Azure Kubernetes Service (AKS), Google Kubernetes Engine (GKE), Oracle Kubernetes Engine (OKE).
\subsubsection{Principali componenti di Kubernetes}
Le applicazioni containerizzate sono dispiegate su componenti detti Pods, questi ultimi sono caricati su uno o più nodi del cluster che possono essere macchine fisiche o virtuali. In ogni nodo sono installati tutti i componenti necessari per controllare i pods e tutti i nodi sono controllati dal nodo master che contiene il layer di gestione del ciclo di vita dei pods.\\
I principali componenti sono quindi:
\begin{itemize}
    \item {
        Nodi: macchine virtuali o fisiche dotate di tutti i componenti per la gestione dei containers. Ogni nodo contiene una container runtime (come CRI o \texttt{containerd}), \texttt{kubelet} (l'agente che gestisce localmente i Pods) ed opzionalmente \texttt{kube-proxy} che implementa una parte del networking.
    }
    \item {
        Pods: la più piccola unità che si può dispiegare su Kubernetes. Ciascun pod contiene uno o più container, tipicamente uno oppure due se strettamente legati ed operano sugli stessi files (sidecar).
    }
    \item {
        Services: in Kubernetes le applicazioni sono distribuite ma forniscono tutte lo stesso \textit{servizio} all'utente o al cluster ed un Service è un modo per unire sotto lo stesso endpoint tutte le applicazioni.\\
        Questo componente permette di gestire la natura effimera dei Pods che possono essere creati e distrutti a runtime quindi l'utente farebbe particolare fatica a tenere traccia dei loro indirizzi.\\
        Fa anche da Load Balancer per distribuire il traffico tra i vari Pods.
    }
    \item {
        Ingress: è un servizio di rete che permette di smistare il traffico in base alle regole indicate dall'utente. Questo permette di erogare diversi servizi provenienti da diversi gruppi di Pods dallo stesso endpoint.\\
        E' componente molto potente in quanto può fare da Load Balancer e come terminazione SSL/TLS; queste caratteristiche riducono la necessità di LBs nel cluster e permette comunicazioni più snelle in quanto il protocollo di sicurezza viene terminato prima dell'ingresso dei pacchetti nel cluster.\\
        E' sempre necessario un Ingress Controller (come AKS, NginX o altri) per applicare le regole di routing.
    }
    \item {
        Autoscaling: HPA ed Autoscalers sono componenti che permettono lo scaling orizzontale (aumentando il numero di Pods) o verticale aumentando il numero di Pods al fine di mantenere entro una certa soglia l'utilizzo delle risorse.\\
        Kubernetes è in grado di comunicare, attraverso plugin come Autoscaler o Karpenter, con i Cloud Providers al fine di eliminare o aggiungere risorse al cluster in base alle necessità ed alle limitazioni. Per esempio potrebbero essere necessari più nodi per i Pod in scheduling oppure rimuovere nodi per abbassare i costi e \textit{consolidare} il cluster.
    }
    \item {
        Storage: la gestione dei dati è fatta dai Volumi. Questi permettono di mantenere dati tra Pods e dopo la loro morte. 
    }
\end{itemize}
%\\
%\\
\textbf{Concludere con la spiegazione di Kubernetes: storage}
%\\
%\\
%
\section{Architettura}
\label{sec:arch}
architettura del cluster, se c'è. Comunque spiegare l'architettura di Kubernetes dentro a OpenStack. 
Fare anche una parte sull'architettura logica dell'applicazione che verrà implementata.
%
\chapter{Implementazione}
Letteralmente
%
\chapter{Testing}
Il Cluster deve andare
%
\chapter{Conclusione e lavori futuri}
Mi piace molto i servizi cloud :)